\chapter{系统需求分析}

需求分析既是系统设计得以推进的起始环节，也是对最终实现是否契合实际业务需求进行检验的重要依据。本章会从功能需求、非功能需求以及角色权限这三个维度出发，对平台的整体需求进行系统梳理，以此为后续的架构设计以及模块划分提供相应依据。

\section{系统功能需求}

平台主要服务于学生、咨询师以及学校管理员这三类用户，不同端的功能边界会依据各自角色职责来进行划分，同时借助统一的后端服务保障数据在各端之间实现一致流转。

\subsection{学生端功能需求}

学生作为平台当中的核心服务群体，其在功能方面的需求主要是围绕"发现资源 $\rightarrow$ 建立连接 $\rightarrow$ 获取支持 $\rightarrow$ 反馈评价"这一较为完整的使用路径来进行展开。

\textbf{个人档案管理。}学生在完成注册之后，可以创建并持续完善个人档案，填写年级、专业以及联系方式等基础信息。相关档案信息会在预约环节自动完成关联，由此使咨询师在受理预约之前，便能够对来访者的基本背景形成初步了解。

\textbf{咨询师信息浏览与筛选。}学生可以对平台内所有处于可预约状态的咨询师进行浏览，查看其专业方向、擅长领域、个人简介以及可预约时段等信息，并在综合比较的基础上自主作出选择。设置该功能的目的在于将更多选择权交还给学生，而不是使其被动接受系统分配。

\textbf{在线预约管理。}学生可以向指定咨询师提交预约申请，并填写预约日期、预约时段以及简要的来访缘由。申请提交之后，预约将进入待确认状态，等待咨询师进行审核。若预约未获通过，或者学生因个人安排临时发生调整，则可以在预约开始之前主动取消预约。与此同时，学生端还提供预约历史列表功能，用于完整展示历次预约在不同阶段的状态变化记录。

\textbf{AI 心理助手对话。}学生可以在任意时间与 AI 心理助手发起交流，以满足情绪倾诉、心理疏导以及心理咨询相关知识查询等方面的需求。对话过程采用流式输出机制，AI 回复将以逐步生成的方式呈现，从而使交互体验更接近实时会话。系统还会自动维护最近 20 条对话记录作为上下文信息，以保障多轮对话在语义层面的连贯性与一致性。

\textbf{公告查阅。}学生可以对学校管理端发布的各类公告信息进行查阅，其中主要包括心理健康活动通知、咨询中心公告以及学校政策说明等内容。系统会按照发布时间倒序进行展示，以便学生能够较为及时地获取最新动态信息。

\textbf{咨询评价提交。}咨询会话结束之后，学生可以对本次咨询体验进行 1—5 分的评分，并提交相应的文字反馈内容。有关评价数据会汇入系统当中，供学校管理端对服务质量开展统计与分析工作。

\subsection{咨询师端功能需求}

咨询师端的功能设计主要围绕"接收预约 $\rightarrow$ 开展会话 $\rightarrow$ 记录存档"这一核心工作流程展开，同时还具备对个人档案进行自主维护的功能。

\textbf{个人档案维护。}咨询师可以在系统当中持续完善并更新个人专业档案，相关内容主要包括姓名、职称、专业方向、擅长问题类型、个人简介以及可预约时段列表等信息。这些内容会直接展示于学生端的咨询师列表页面，并成为学生在进行咨询师选择时所依据的重要参考信息。

\textbf{预约审核管理。}咨询师可以集中查看当前全部处于待处理状态的预约申请，并在阅读来访者提交的简要说明之后，对该预约作出接受或者拒绝的审核决定。若选择拒绝，还可以进一步补充相应的原因说明，并由系统自动向学生发送结果通知。对于已经完成确认的预约，系统会将其同步纳入咨询师的工作安排视图当中，以便后续对咨询时间开展统筹与管理。

\textbf{会话记录管理。}当预约的咨询时间到达后，咨询师可以在系统当中创建新的会话记录，并将其与对应的预约信息完成关联。待会话结束之后，咨询师可执行"结束会话"操作，系统会同步记录本次会话时长，并对相关状态进行更新。历史会话列表支持按照时间顺序进行查看，以便咨询师持续跟踪长期个案的进展情况。

\textbf{案例笔记撰写。}针对每一次咨询会话，咨询师均可在系统内部撰写私密案例笔记，用于记录来访者主诉、会谈重点、评估结论以及后续干预计划。该类笔记仅向咨询师本人开放，不会在学生端和学校管理端显示，从而进一步保障来访者的隐私安全。与此同时，笔记内容支持多次补充与更新，以满足跨多次会话的长期咨询记录需求。

\subsection{学校管理端功能需求}

学校管理端主要面向心理中心管理人员，其功能定位主要集中在资源调配、服务监管以及信息发布等方面。

\textbf{咨询师账号管理。}管理员可以创建咨询师账号，完成初始密码设置，并将账号与对应的咨询师档案进行绑定。对于已经离职或者阶段性暂停服务的咨询师，管理员可以执行禁用操作，使其不再在学生端显示，同时确保相关历史数据得到完整保留。

\textbf{平台数据统计。}管理端具备平台级数据概览与可视化统计功能。数据概览卡片集中呈现在读学生人数、在职咨询师人数、累计预约总量、已完成会话数量及平均满意度等核心指标；图表统计模块进一步以环形饼图展示各咨询师预约量占比、以环形饼图展示咨询话题分布（学业压力、情感问题、人际关系、家庭问题、就业焦虑、心理健康、其他七类），并以柱状图呈现预约量最多的前十名学生，从而为管理层开展服务质量评估与资源调配决策提供多维度数据支撑。学生在发起预约时需选择所属话题类别，该结构化字段作为话题统计的数据来源。

\textbf{公告发布与管理。}管理员可以在系统中完成公告标题以及正文内容的编辑并提交，使公告能够及时在学生端与咨询师端进行展示。系统同时支持对已不再适用的历史公告按需进行删除处理。

\subsection{功能需求汇总}

图~\ref{fig:2-func}~以功能分解图的方式，对三类角色各自可开展的功能操作进行了较为直观的展示。学生端围绕"发现—预约—对话—评价"的使用流程来进行设计，共设置 6 项功能；咨询师端以"接收—记录—归档"的工作流程为核心，共设置 4 项功能；学校管理端则主要聚焦于平台运营管理方面，共设置 3 项功能。

\begin{figure}[H]
    \centering
    \begin{tikzpicture}[
        role/.style={draw, rounded corners=4pt, minimum width=3.2cm,
                     minimum height=0.75cm, text centered,
                     font=\heiti\small, inner sep=4pt},
        func/.style={draw, rounded corners=3pt, minimum width=3.2cm,
                     minimum height=0.6cm, text width=3.0cm,
                     text centered, font=\small, inner sep=3pt},
        arr/.style={->, thick, gray!70}
    ]
    %--- 学生端（左列）---
    \node[role, fill=cyan!20, draw=cyan!60]
        (SR) at (0,0) {学生端};
    \node[func, fill=cyan!5,  below=0.35cm of SR] (S1) {个人档案管理};
    \node[func, fill=cyan!5,  below=0.18cm of S1] (S2) {咨询师浏览与筛选};
    \node[func, fill=cyan!5,  below=0.18cm of S2] (S3) {在线预约管理};
    \node[func, fill=cyan!5,  below=0.18cm of S3] (S4) {AI 心理助手对话};
    \node[func, fill=cyan!5,  below=0.18cm of S4] (S5) {公告查阅};
    \node[func, fill=cyan!5,  below=0.18cm of S5] (S6) {咨询评价提交};
    \draw[arr](SR)--(S1); \draw[arr](S1)--(S2); \draw[arr](S2)--(S3);
    \draw[arr](S3)--(S4); \draw[arr](S4)--(S5); \draw[arr](S5)--(S6);

    %--- 咨询师端（中列）---
    \node[role, fill=orange!20, draw=orange!60]
        (CR) at (5.5,0) {咨询师端};
    \node[func, fill=orange!5, below=0.35cm of CR] (C1) {个人档案维护};
    \node[func, fill=orange!5, below=0.18cm of C1] (C2) {预约审核管理};
    \node[func, fill=orange!5, below=0.18cm of C2] (C3) {会话记录管理};
    \node[func, fill=orange!5, below=0.18cm of C3] (C4) {案例笔记撰写};
    \draw[arr](CR)--(C1); \draw[arr](C1)--(C2);
    \draw[arr](C2)--(C3); \draw[arr](C3)--(C4);

    %--- 学校管理端（右列）---
    \node[role, fill=green!20, draw=green!60]
        (AR) at (11,0) {学校管理端};
    \node[func, fill=green!5,  below=0.35cm of AR] (A1) {咨询师账号管理};
    \node[func, fill=green!5,  below=0.18cm of A1] (A2) {平台数据统计};
    \node[func, fill=green!5,  below=0.18cm of A2] (A3) {公告发布与管理};
    \draw[arr](AR)--(A1); \draw[arr](A1)--(A2); \draw[arr](A2)--(A3);
    \end{tikzpicture}
    \caption{系统各角色功能分解图}
    \label{fig:2-func}
\end{figure}

表~\ref{tab:2-1}~进一步对三端功能需求进行详细汇总。

\begin{table}[H]
    \centering
    \caption{系统功能需求汇总}
    \label{tab:2-1}
    \begin{tabularx}{\linewidth}{lXl}
        \toprule
        功能模块 & 子功能 & 涉及角色 \\
        \midrule
        认证模块 & 注册、登录、Token 刷新 & 全部 \\
        个人档案 & 查看档案、更新档案 & 学生、咨询师 \\
        咨询师浏览 & 获取可用咨询师列表 & 学生 \\
        预约管理 & 发起预约、取消预约、查看预约列表 & 学生 \\
        预约审核 & 确认预约、拒绝预约、查看预约列表 & 咨询师 \\
        会话管理 & 创建会话、结束会话、查看会话列表 & 咨询师 \\
        案例笔记 & 新增笔记、查看笔记 & 咨询师 \\
        AI 心理助手 & 发起对话、查看历史对话 & 学生 \\
        咨询评价 & 提交评分与文字反馈 & 学生 \\
        公告管理 & 发布、删除、查看公告列表 & 管理员（发布）、学生/咨询师（查看） \\
        咨询师管理 & 添加账号、启用/禁用账号 & 管理员 \\
        数据统计 & 平台总览数据查询 & 管理员 \\
        \bottomrule
    \end{tabularx}
\end{table}


\section{系统非功能需求}

功能需求主要用于说明系统需要完成哪些工作，非功能需求则进一步界定系统在完成这些工作时的质量表现。对于面向高校应用场景构建的心理咨询平台来说，以下几类非功能属性属于系统设计与运行过程中较为关键的约束条件。

\subsection{性能需求}

\textbf{响应时间。}对于普通 API 请求（比如数据查询以及表单提交等操作），系统在正常网络环境下应于 500 毫秒内完成响应；对于涉及数据库关联查询的复杂接口，则需要把响应时间控制在 1 秒以内。AI 流式对话接口会受到外部大模型 API 网络延迟的影响，因此其首字节时间（TTFB）应控制在 3 秒以内，后续流式输出应保持连续，并避免出现明显卡顿。

\textbf{并发能力。}平台在建设初期以单校部署作为基本定位，预期可支持 100 名学生同时在线使用，覆盖普通页面浏览以及 AI 对话等并发使用场景。Node.js 依托事件循环的非阻塞 I/O 模型，较为适宜对这类高并发、I/O 密集型的业务请求进行处理。

\textbf{数据库查询效率。}应针对 \texttt{user\_id}、\texttt{appointment\_status} 以及 \texttt{counselor\_id} 等高频查询字段进行索引配置，从而保证在数据量持续增长之后，关键查询的响应时间仍能保持稳定，不会出现明显下降。

\subsection{安全性需求}

\textbf{身份认证。}系统借助 JSON Web Token 来实现无状态的认证机制，同时把 Token 的有效期统一设置为 7 天。每次发起请求时，都需要在 Authorization 请求头中携带有效的 Token，只有在服务端完成校验之后，才可以访问受保护资源；如果 Token 已经过期或存在伪造情况，接口则会返回 401 状态码，并给出重新登录提示。

\textbf{密码安全。}用户密码在注册阶段会先经过 cost factor = 10 的 bcrypt 算法进行哈希处理，之后再写入数据库，系统在任何处理环节中都不会对明文密码进行存储或传输。登录验证时则借助 bcrypt 的比较函数来完成校验，因此即便数据库发生泄露，攻击者也难以凭借查表方式还原出原始密码。

\textbf{权限控制。}三端接口均借助角色守卫中间件来进行细粒度的访问控制，其中 \texttt{student} 角色不能访问咨询师端以及管理端接口，反向的跨角色访问同样会被系统禁止。对于任何角色发起的越权请求，系统都会返回 403 状态码，并在拒绝访问的同时对异常情况进行记录。

\textbf{SQL 注入防护。}系统在开展数据库交互相关工作时统一借助 Sequelize ORM 来完成，所有查询参数都会按照参数化的方式进行传递，从框架机制这一层面对 SQL 注入风险进行阻断，同时也避免手工去编写 SQL 语句。

\subsection{可用性需求}

\textbf{界面友好性。}三端界面借助不同主题色来进行区分，学生端选用蓝绿色系，咨询师端选用暖橙色系，管理端选用深蓝色系，以便帮助用户较快形成清晰的视觉归属。预约、对话以及审核等核心操作流程应尽量控制在 3 次点击以内，这样可以减少因操作路径过长而带来的认知负担。

\textbf{操作反馈及时性。}对于表单提交、状态更新等操作，系统应在处理完成后及时给出成功或失败的结果提示，从而避免用户因缺少反馈信息而出现重复提交。AI 对话则借助带有打字机效果的流式输出方式进行呈现，使用户能够感知系统仍处于处理中，进而减少等待过程中的不确定感。

\textbf{兼容性。}系统把现代浏览器（Chrome 90+、Edge 90+以及 Firefox 88+）作为目标运行环境，无需额外安装客户端软件，从而能够适应高校场景下设备类型多样的使用需求。

\subsection{可维护性与可扩展性}

\textbf{分层架构。}后端整体严格遵循 Controller $\rightarrow$ Service $\rightarrow$ Model 的三层分离式设计，把业务逻辑统一放在 Service 层当中，而 Controller 仅负责对请求进行解析，并对响应内容进行组装。这样的结构可以使单个模块在发生改动时，不会对其他模块带来连锁影响。

\textbf{模块化路由。}系统把不同角色对应的接口分别划分到独立的路由文件中进行管理，因此当需要新增业务端点时，只需要在对应文件内补充相关路由定义，而不必改动核心应用文件。

\textbf{数据库迁移管理。}数据库结构的变更统一借助按编号递增追加的 migration 文件来进行管理，例如 \texttt{002\_add\_xxx.sql}；每个 migration 均运用 \texttt{IF NOT EXISTS} 等幂等语法，从而确保脚本可以被重复执行，并为后续功能扩展保留规范化的数据库演进路径。


\section{系统角色与权限}

\subsection{三类角色定义}

系统共定义三类用户角色，并在用户注册或账号创建阶段将角色信息写入 JWT Payload，后续每一次请求均由服务端统一完成校验。

\textbf{学生（student）。}系统的主要服务对象，账号通常借助用户自助注册的方式来获取。该角色的权限范围会受到严格限制，仅可访问与本人相关的数据内容，包括个人预约信息、个人对话历史以及个人评价记录，不能查看其他用户的相关信息。

\textbf{咨询师（counselor）。}平台中承担心理咨询服务工作的专业人员，账号由学校管理员统一进行创建并完成授权。该角色仅可访问与自身工作直接相关的预约、会话以及笔记等数据，并可对归属本人的预约执行确认或拒绝处理，但不能查看其他咨询师的会话记录，也无法使用学校管理端的相关功能。

\textbf{学校管理员（school）。}平台开展运营管理工作的核心角色，一般会由心理中心主任或者相关管理人员来担任。管理员账号会在系统初始化的阶段完成预先设置，拥有对咨询师账号进行管理、对整个平台数据开展统计以及发布公告的权限，但不会参与到具体的咨询会话内容当中，以保护咨访关系的私密性。

\subsection{权限矩阵}

表~\ref{tab:2-2}~从资源维度展示三类角色的访问权限矩阵。

\begin{table}[H]
    \centering
    \caption{系统权限矩阵}
    \label{tab:2-2}
    \begin{tabularx}{\linewidth}{Xccc}
        \toprule
        资源 / 操作 & 学生 & 咨询师 & 管理员 \\
        \midrule
        个人档案（本人） & 读 / 写 & 读 / 写 & — \\
        咨询师列表 & 只读 & — & 读 / 写 \\
        预约（本人发起） & 发起 / 取消 / 查看 & — & — \\
        预约（本人接收） & — & 确认 / 拒绝 / 查看 & — \\
        会话记录 & — & 创建 / 结束 / 查看 & — \\
        案例笔记 & — & 读 / 写（仅本人） & — \\
        AI 对话 & 发起 / 查看（本人） & — & — \\
        咨询评价 & 提交 & — & 只读（统计） \\
        公告 & 只读 & 只读 & 读 / 写 / 删 \\
        咨询师账号管理 & — & — & 创建 / 启用 / 禁用 \\
        平台统计数据 & — & — & 只读 \\
        \bottomrule
    \end{tabularx}
\end{table}


\section*{本章小结}

本章从功能需求、非功能需求以及角色权限这三个维度出发，对平台整体的设计要求进行了较为系统的梳理。在功能需求方面，已对学生端、咨询师端以及学校管理端各自对应的功能边界和核心使用路径进行了明确梳理，并借助汇总表把 12 个功能模块的完整规格集中展示出来；在非功能需求方面，分别从性能、安全、可用性以及可维护性这四个方向出发，设置了量化或者定性的约束条件；在角色权限方面，对三类用户角色的职责范围进行了界定，并以权限矩阵的形式进一步明确了各角色对不同资源所拥有的访问权限。这些分析结果会为下一章系统架构与数据库设计提供直接的依据。
